\documentclass[UTF8,zihao=-4]{ctexart}
\usepackage[a4paper,margin=2.2cm]{geometry}
\usepackage{fontspec}
\usepackage{tabularx}
\usepackage{array}
\usepackage{ragged2e}
\usepackage{fancyhdr}
\usepackage{hyperref}
\setCJKmainfont{Noto Serif SC}
\setCJKsansfont{Noto Sans SC}
\setmainfont{Noto Serif SC}
\setlength{\parindent}{0pt}
\setlength{\parskip}{0.45em}
\pagestyle{fancy}
\fancyhf{}
\fancyhead[L]{商务智能}
\fancyfoot[C]{\thepage}
\hypersetup{colorlinks=true,linkcolor=black,urlcolor=blue}
\begin{document}
\title{10、建立数据集市原因}
\author{}
\date{}
\maketitle

企业级数据仓库通常规模较大，虽然提供统一数据模式、统一数据表示和统一访问接口，但终端用户在具体部门或应用场景下仍然需要面向局部主题的快速分析环境。网络环境中即使存在多个数据源，最终用户仍可能得不到可用信息。早期数据仓库主要强调多数据源集成，为用户访问多个数据源提供统一视图。因此，可以按照部门、个人、子主题或业务区域建立数据集市，以满足特定商务单元、商务程序或商务应用的需求。\par

\end{document}
